\documentclass[11pt]{article}
\usepackage[T1]{fontenc}
\usepackage{amsmath,amssymb,amsthm,lmodern}
\usepackage[margin=1in]{geometry}
\usepackage{hyperref}
\hypersetup{hidelinks,pdftitle={EGNO Theorem 4.4.1: absence of self-extensions of the unit},pdfauthor={Pavel Etingof; Shlomo Gelaki; Dmitri Nikshych; Victor Ostrik}}
\DeclareMathOperator{\Hom}{Hom}
\DeclareMathOperator{\End}{End}
\DeclareMathOperator{\Ext}{Ext}
\DeclareMathOperator{\Rep}{Rep}
\newcommand{\Cat}{\mathcal C}
\newcommand{\unit}{\mathbf 1}
\newcommand{\BookVec}{\operatorname{Vec}}
\newtheorem*{primarythm}{Theorem 4.4.1}
\title{Tensor Categories\\Theorem 4.4.1: Absence of Self-Extensions of the Unit}
\author{Pavel Etingof \and Shlomo Gelaki \and Dmitri Nikshych \and Victor Ostrik}
\date{Private reader of the 2015 AMS author final copy}
\begin{document}
\maketitle

% BEGIN CONTEXT EGNO:4.4.1
\noindent\textbf{Fixed source and scope.}
\emph{Tensor Categories}, Mathematical Surveys and Monographs 205,
American Mathematical Society, 2015, author final copy.
\par\noindent\url{https://math.mit.edu/~etingof/egnobookfinal.pdf}.
\par\noindent
The retained PDF has 362 physical pages; its metadata record creation on
May 13, 2015 and modification on July 29, 2015. The primary is the
statement on physical 87 (printed 70) and its complete proof on physical
88 (printed 71). Corollary 4.4.2 and Remark 4.4.3 are uncounted context.

Chapter 4 opens on physical 82 (printed 65) with ``Let $k$ be an
algebraically closed field.'' Section 4.4 adds characteristic zero.
The theorem assumes a finite \emph{ring} category and a simple unit;
it does not assume left or right duals, rigidity, symmetry,
semisimplicity, or a fiber functor. The original ``(right) action'' and
``homomorphism'' wording, ``etc.'' and ``We can now proceed further''
are retained. The passage to higher tensor powers and the final
finite-character sentence are not expanded into a new proof.

\noindent\textbf{Source notice.}
The source footer states: ``Author's final version made available with
permission of the publisher, American Mathematical Society. License or
copyright restrictions may apply to redistribution; see
\url{http://www.ams.org/publications/ebooks/terms}.'' This reader remains
private. The separate February 14, 2026 corrections do not silently
replace the 2015 text.

\subsection*{Definition context (uncounted)}
Definition 4.2.3 (physical 83, printed 66) reads: a \emph{multiring
category} over $k$ is a locally finite $k$-linear abelian monoidal
category $\Cat$ with bilinear and biexact tensor product. If in addition
$\End_{\Cat}(\unit)=k$, we will call $\Cat$ a \emph{ring category}.
No duals occur in this definition.

Definition 1.8.5 (physical 26, printed 9) says that a $k$-linear abelian
category $\Cat$ is \emph{finite} if it is equivalent to the category
$A\text{-mod}$ of finite dimensional modules over a finite dimensional
$k$-algebra $A$. Definition 1.8.6 gives the equivalent conditions:
finite dimensional Hom spaces, finite length of every object, enough
projectives (every simple object has a projective cover), and finitely
many isomorphism classes of simple objects. The first two conditions
alone define \emph{locally finite} (Definition 1.8.1). The equivalence
of these two definitions, discussed on physical 27 (printed 10), is an
explicit prior input; its proof is not reprinted.

Definition 1.6.5 (physical 23, printed 6) calls $P$ projective if
$\Hom_{\Cat}(P,-)$ is exact. Under the finite-length condition,
Definition 1.6.6 calls $p:P(X)\to X$ a \emph{projective cover} if
$P(X)$ is projective, $p$ is an epimorphism, and for every epimorphism
$g:P\to X$ from a projective object there is an epimorphism
$h:P\to P(X)$ such that $ph=g$. Such a cover, when it exists, is
unique up to a non-unique isomorphism.

Definition 1.4.2 (physical 21, printed 4) defines morphisms between
short exact sequences $0\to X\to Z\to Y\to0$ and
$0\to X\to Z'\to Y\to0$ as maps $Z\to Z'$ restricting to the
identity on $X$ and inducing the identity on $Y$. The set of these
sequences up to isomorphism is $\Ext^1(Y,X)$, the set of extensions
of $Y$ by $X$. Its addition and identification with derived Ext are
standard prior background (Section 1.4 and physical 24, printed 7);
their proofs are not reprinted. The notation $A^*$ in the proof below
means the vector-space dual of the finite dimensional algebra $A$.
The ordinary symbol $\BookVec$ denotes finite dimensional $k$-vector spaces.
% END CONTEXT EGNO:4.4.1

% BEGIN PRIMARY STATEMENT EGNO:4.4.1
\subsection*{4.4. Absence of self-extensions of the unit object}
Assume that $k$ has characteristic $0$.
\begin{primarythm}
Let $\Cat$ be a finite ring category over $k$ with simple object $\unit$.
Then $\Ext^1(\unit,\unit)=0$.
\end{primarythm}
% END PRIMARY STATEMENT EGNO:4.4.1

% BEGIN PRIMARY PROOF EGNO:4.4.1
\begin{proof}
Assume the contrary, and suppose that $V$ is a nontrivial extension of
$\unit$ by itself. Let $P$ be the projective cover of $\unit$. Then
$\Hom(P,V)$ is a 2-dimensional space, with a filtration induced by the
filtration on $V$, and both quotients naturally isomorphic to
$E:=\Hom(P,\unit)$. Let $v_0,v_1$ be a basis of $\Hom(P,V)$ compatible
to the filtration, i.e., $v_0$ spans the 1-dimensional subspace defined
by the filtration. Let $A=\End(P)$ (this is a finite dimensional
algebra). Let $\varepsilon:A\to k$ be the character defined by the
(right) action of $A$ on $E$. Then the matrix of $a\in A$ in the basis
$v_0,v_1$ has the form
\begin{equation}
[a]_1=\begin{pmatrix}
\varepsilon(a)&\chi_1(a)\\
0&\varepsilon(a)
\end{pmatrix}
\tag{4.5}
\end{equation}
where $\chi_1\in A^*$ is nonzero. Since $a\mapsto[a]_1$ is a
homomorphism, $\chi_1$ is a derivation:
$\chi_1(xy)=\chi_1(x)\varepsilon(y)+\varepsilon(x)\chi_1(y)$.

Now consider the representation $V\otimes V$. Using the exactness of
the tensor product, we see that the space $\Hom(P,V\otimes V)$ is
4-dimensional, and has a 3-step filtration, with successive quotients
$E,E\oplus E,E$, and basis $v_{00};v_{01},v_{10};v_{11}$, consistent
with this filtration. The matrix of $a\in\End(P)$ in this basis is
\begin{equation}
[a]_2=\begin{pmatrix}
\varepsilon(a)&\chi_1(a)&\chi_1(a)&\chi_2(a)\\
0&\varepsilon(a)&0&\chi_1(a)\\
0&0&\varepsilon(a)&\chi_1(a)\\
0&0&0&\varepsilon(a)
\end{pmatrix}.
\tag{4.6}
\end{equation}
Since $a\mapsto[a]_2$ is a homomorphism, we find
\[
\chi_2(ab)=\varepsilon(a)\chi_2(b)+\chi_2(a)\varepsilon(b)
+2\chi_1(a)\chi_1(b),\qquad a,b\in A.
\]
We can now proceed further (i.e., consider $V\otimes V\otimes V$
etc.) and define for every positive $n$, a linear function
$\chi_n\in A^*$ which satisfies the equation
\[
\chi_n(ab)=\sum_{j=0}^n\binom nj\chi_j(a)\chi_{n-j}(b),
\]
where $\chi_0=\varepsilon$. Thus for any $s\in k$, we can define
$\varphi_s:A\to k((t))$ by
$\varphi_s(a)=\sum_{m\ge0}\chi_m(a)s^mt^m/m!$ (this is where we use
that $k$ has characteristic $0$) and we find that $\varphi_s$ is a
family of pairwise distinct homomorphisms. This is a contradiction,
as $A$ is a finite dimensional algebra and can have only finitely many
1-dimensional representations over any extension field. We are done.
\end{proof}
% END PRIMARY PROOF EGNO:4.4.1

% BEGIN SUPPORT EGNO:4.4.1 CONSEQUENCE
\subsection*{Consequence and scope contrast (uncounted)}
\noindent\emph{Physical 88, printed 71; complete printed passages.}
\paragraph{Corollary 4.4.2.}
If a finite ring category $\Cat$ over a field of characteristic zero
has a unique simple object $\unit$, then $\Cat$ is equivalent to the
category $\BookVec$.

\paragraph{Remark 4.4.3.}
The conclusion of Corollary 4.4.2 does not hold when the ground field
has a positive characteristic. We give a counterexample: the category
$\Rep(\mathbb Z/p\mathbb Z)$ over a field of characteristic $p$.
% END SUPPORT EGNO:4.4.1 CONSEQUENCE

\subsection*{Separate 2026 corrections and source credit}
The six-page author correction document is dated February 14, 2026:
\url{https://math.mit.edu/~etingof/tensbookcor.pdf}.
No entry in its Chapter 4 list changes Theorem 4.4.1 or its displayed
proof. The physical-page-3 entries change the reference in the proof
of Proposition 4.2.8, interchange ``monomorphisms'' and
``epimorphisms'' in Corollary 4.3.9, and require finiteness in the
paragraph before Example 4.5.5. Those passages are not imported as
hypotheses or silently substituted into this proof. The primary's
finite-ring and characteristic-zero hypotheses already appear in
the 2015 source.

The four book authors supply the exposition transcribed here.
Note 4.4 on physical 105 (printed 88) states: ``The results on absence
of self-extensions are taken from [EtO1] (see Theorem 2.17).''
The cached book bibliography on physical 345 (printed 328) gives:
\begin{quote}
[EtO1] P. Etingof and V. Ostrik. \emph{Finite tensor categories}.
Moscow Math. Journal 4 (2004), 627--654.
\end{quote}
This is the book's historical credit, not a claim that the separate
paper was obtained or that its other results are included. No new
derivation of the higher-$\chi_n$ construction, finite-character
assertion, or other author shorthand is supplied.
\end{document}
